\documentclass[../../main.tex]{subfiles}
\graphicspath{{../../image/}{../../figure/}}
\begin{document}

\section{弱解的非唯一性}\label{section:ch04-弱解的非唯一性}

把解的概念从古典解推广到弱解引入了一个新的困难：非唯一性。为看到这一点，考虑\eqref{align:ch04-Burgers}以初值
\begin{align}
u(x,0)=\begin{cases}-1,&x<0\\ 1,&x>0\end{cases}
\label{align:ch04-Riemann初值}
\end{align}
的 Cauchy 问题。这是著名的 Riemann 问题的一个例子，将在第 IX 章中详加讨论。问题\eqref{align:ch04-Burgers}、\eqref{align:ch04-Riemann初值}有无穷多个弱解，包括下面的解族
\begin{align}
u_{\alpha}(x,t)=
\begin{cases}
-1,&-\infty<x\leqslant-t\\
\frac{x}{t},&-t<x\leqslant-\alpha t\\
-\alpha,&-\alpha t<x\leqslant 0\\
\alpha,&0<x\leqslant\alpha t\\
\frac{x}{t},&\alpha t<x\leqslant t\\
1,&t<x<\infty,
\end{cases}
\label{align:ch04-非唯一解族}
\end{align}
其中任一 $\alpha\in[0,1]$。事实上，只要 $x/t\notin\{0,\pm\alpha,\pm 1\}$，$u_{\alpha}(x,t)$ 就在古典意义下满足\eqref{align:ch04-Burgers}。直线 $x/t=\pm 1$（对 $\alpha\in[0,1]$）与直线 $x/t=\pm\alpha$（对 $\alpha\in(0,1)$）只是弱间断面，$u_{\alpha}$ 跨越它们连续。最后，当 $\alpha\neq 0$ 时，直线 $x=0$ 是静止激波面，跨越它时 Rankine--Hugoniot 跳跃条件\eqref{align:ch04-RH}成立。

为解决非唯一性问题，须对弱解附加额外的限制，即\emph{可容性条件}。一开始，合理的可容性判据至少应满足以下要求中的若干条：
\begin{enumerate}
\item\label{item:ch04-可容-a} 它们应由物理决定，或至少为物理所启示。
\item\label{item:ch04-可容-b} 它们应与其他已建立的可容性条件相容。
\item\label{item:ch04-可容-c} 它们应足够宽泛，以容许可容解的存在；同时又足够狭窄，以排除赝解。理想情况下，应能从中挑选出唯一的可容解。
\end{enumerate}

双曲守恒律弱解的可容性问题将是本书的中心主题，需要冗长的讨论，这将在接下来两节开始，并在第 VIII、IX 章达到高潮。特别地，我们将会看到，$u_{0}(x,t)$ 是本节所考虑的简单问题\eqref{align:ch04-Burgers}、\eqref{align:ch04-Riemann初值}的唯一可容解。

\end{document}
